\section{Proof of Theorem~\ref*{thm:sharper}(1): a counterexample below six}\label{sec:subsixproof}

This appendix gives the counterexample with both factor capacities
below six, proving \cref{thm:sharper}(1). We specify the graphs, codes
and exact matrix certificates that establish the required capacity
bounds, together with procedures for checking them.

The proof in this appendix combines analytical arguments with
computer-assisted verification of explicit graph data, codes and matrix
certificates. OpenAI Codex assisted with exploring candidate constructions
and certificates and with developing and reviewing the verification code.
The finite certificate checks use exact arithmetic, and the
product-capacity bound follows from the analytical spectral argument.

\subsection{Statement}

We use the conventions of Supplementary Apps.~\ref{sec:notation}--\ref{sec:spectrum}.
A length-$n$ code in $X$ is a clique in $X^{\OR n}$, and
$\C(X)=\sup_{n\ge1}\omega(X^{\OR n})^{1/n}$.

Set
\[
 b=\frac{32437757}{6250000}=5.19004112,\quad
 q=\frac{12991403}{2500000}=5.1965612,\quad
 L=\frac{5193957}{1000000}=5.193957.
\]
\begin{proposition}\label{subsix:thm:main}
For the graphs \(G,H\) defined in Subsec.~\ref{subsix:sec:graphs},
\[
 \C(G\times H)<L<\sqrt{27}\leq\min\{\C(G),\C(H)\},
\]
and, separately,
\[
 \C(G)\leq\frac{16}{3}<6,\qquad \C(H)\leq q<6.
\]
Consequently the Shannon OR-capacity Hedetniemi equality fails.
The same pair witnesses
\(\C(G\times H)<r<\min\{\C(G),\C(H)\}\) for every real
\(r\in[L,\sqrt{27})\).
\end{proposition}
This proves \cref{thm:sharper}(1).
The smaller factor capacity lies in $[\sqrt{27},q]$;
its exact value is unknown. The construction establishes this local
interval, with no claim that its lower endpoint is optimal or that
failure extends to every larger threshold.

\subsection{The graphs and their lower-bound codes}\label{subsix:sec:graphs}

\subsubsection{The 54-vertex graph}

Use the following ordered list of vectors in \(\F_3^5\), with entries
interpreted modulo three.
\begin{center}
\small
\begin{tabular}{r@{\quad}l r@{\quad}l r@{\quad}l}
\toprule
\(i\)&\(w_i\)&\(i\)&\(w_i\)&\(i\)&\(w_i\)\\
\midrule
0&(0,0,0,0,1) & 9&(0,1,2,0,0) & 18&(1,1,1,2,1)\\
1&(0,0,0,1,1) & 10&(0,1,2,2,1) & 19&(1,1,1,2,2)\\
2&(0,0,0,1,2) & 11&(0,1,2,2,2) & 20&(1,1,2,1,1)\\
3&(0,0,1,2,0) & 12&(1,0,0,0,0) & 21&(1,1,2,2,0)\\
4&(0,1,0,1,0) & 13&(1,0,1,0,0) & 22&(1,2,1,1,0)\\
5&(0,1,0,2,0) & 14&(1,0,1,1,2) & 23&(1,2,1,2,1)\\
6&(0,1,1,0,0) & 15&(1,0,2,0,0) & 24&(1,2,1,2,2)\\
7&(0,1,1,1,1) & 16&(1,0,2,2,2) & 25&(1,2,2,1,1)\\
8&(0,1,1,1,2) & 17&(1,1,1,1,0) & 26&(1,2,2,2,0)\\
\bottomrule
\end{tabular}
\end{center}
Let \(B\) have vertex set \(\{0,\ldots,26\}\) and adjacency
\[
 i\sim_Bj\quad\Longleftrightarrow\quad
 i\ne j\ \text{and}\ w_i\cdot w_j=0\text{ in }\F_3.
\]
Let \(D=\overline B\), the simple complement on the same vertex set.
Define
\[
 G=B\sqcup D.
\]
For explicit labels, the \(B\)-copy uses \(0,\ldots,26\) and the \(D\)-copy
uses \(27,\ldots,53\), with \(27+i\) representing vertex \(i\) of \(D\).
There are no edges between the copies.

The \(27\) words
\[
 \mathcal K_G=\{(i,27+i):0\leq i\leq26\}\subseteq V(G)^2
\]
form a code because, for distinct \(i,j\), exactly one of \(i\sim_Bj\) and
\(i\sim_Dj\) holds. Therefore
\begin{equation}\label{subsix:eq:Glower}
 \C(G)\geq\omega(G\OR G)^{1/2}\geq\sqrt{27}.
\end{equation}

\subsubsection{The 27-vertex graph}

Let \(V(H)=\mathbb Z_{29}\setminus\{1,17\}\), ordered as
\(0,2,3,\ldots,16,18,19,\ldots,28\). The following table defines the graph completely
by listing, in row \(v\), all its neighbours larger than \(v\).
Each listed pair is an undirected edge, and there are no loops.
The residue labels specify the vertices of \(H\); adjacency is defined
entirely by the table.
\begin{longtable}{r@{\quad}p{0.84\textwidth}}
\toprule
\(v\)&Larger neighbours in \(H\)\\
\midrule
\endfirsthead
\toprule
\(v\)&Larger neighbours in \(H\), continued\\
\midrule
\endhead
0&\(2,\,4,\,6,\,7,\,8,\,12,\,13,\,16,\,21,\,22,\,23,\,25,\,27\)\\
2&\(4,\,5,\,6,\,7,\,8,\,9,\,10,\,12,\,15,\,22,\,23\)\\
3&\(8,\,9,\,11,\,14,\,15,\,20,\,22,\,23,\,26,\,28\)\\
4&\(5,\,6,\,7,\,8,\,10,\,13,\,15,\,18,\,20,\,22,\,23,\,25,\,28\)\\
5&\(6,\,7,\,9,\,11,\,12,\,13,\,16,\,18,\,19,\,21,\,22,\,24\)\\
6&\(8,\,9,\,10,\,15,\,18,\,20,\,25,\,26,\,27\)\\
7&\(10,\,11,\,13,\,14,\,18,\,19,\,20,\,24,\,25,\,26,\,27,\,28\)\\
8&\(9,\,10,\,12,\,16,\,24\)\\
9&\(10,\,11,\,12,\,13,\,15,\,16,\,20,\,22,\,23,\,25,\,26\)\\
10&\(11,\,13,\,22,\,24,\,28\)\\
11&\(12,\,13,\,14,\,15,\,22,\,23,\,24,\,25\)\\
12&\(13,\,14,\,15,\,16,\,18,\,20,\,21,\,24,\,27,\,28\)\\
13&\(15,\,16,\,20,\,21,\,24\)\\
14&\(15,\,16,\,18,\,20,\,23,\,25,\,26,\,27,\,28\)\\
15&\(18,\,22,\,26,\,28\)\\
16&\(18,\,19,\,20,\,22,\,24,\,25\)\\
18&\(19,\,20,\,22,\,24,\,26\)\\
19&\(20,\,21,\,22,\,23,\,25,\,27,\,28\)\\
20&\(21,\,23,\,24,\,26,\,27\)\\
21&\(23,\,25,\,26,\,27\)\\
22&\(24,\,25,\,27,\,28\)\\
23&\(24,\,25,\,27\)\\
24&\(25,\,27\)\\
25&\(27,\,28\)\\
26&\(28\)\\
27&\(\varnothing\)\\
28&\(\varnothing\)\\
\bottomrule
\end{longtable}
This graph has \(176\) edges. The \(27\) words
\[
 \mathcal K_H=\{(i,12i\bmod29):i\in\mathbb Z_{29}\setminus\{1,17,28\}\}
 \cup\{(28,12)\}
\]
are all in \(V(H)^2\), since \(12i=1\) and \(12i=17\) modulo \(29\)
have the unique solutions \(i=17\) and \(i=28\), respectively.
Their first coordinates are distinct, and the table verifies
\[
 (a,b),(c,d)\in\mathcal K_H,\quad(a,b)\ne(c,d)
 \ \Longrightarrow\
 a\sim_Hc\ \text{or}\ b\sim_Hd.
\]
This is a finite check of \(351\) unordered pairs, also performed by
the exact verifier described in Subsec.~\ref{subsix:app:verification}.
Hence
\begin{equation}\label{subsix:eq:Hlower}
 \C(H)\geq\sqrt{27}.
\end{equation}

\subsection{The finite-field certificates}\label{subsix:sec:rank}

A block representation over $\F_3$ of a graph \(X\), with local block size \(d\)
and ambient dimension \(R\), consists of matrices
\(P_v\in\F_3^{d\times R}\), \(Q_v\in\F_3^{R\times d}\) such that
\[
 P_vQ_v=I_d,\qquad P_vQ_w=0\quad\text{whenever }v\sim_Xw.
\]
Both orientations of an edge are included. The resulting block matrix
$(P_vQ_w)_{v,w}$ need not be symmetric and has rank at most $R$.
By the definition of fractional Haemers' bound in
\suppref{sec:spectrum} and \cref{lem:haemers}, these factors certify
\[
 h_3(X):=\phif{\F_3}(X)\le R/d.
\]
Both $h_3$ and $\overline\vartheta(X)=\vartheta(\overline X)$ are OR
spectral points, as recalled in \cref{sec:spectral-argument} and
\cref{lem:haemers}.

\subsubsection{A rank-five scalar representation of B}

Every listed vector satisfies \(w_i\cdot w_i\ne0\) in \(\F_3\).
Define the row and column factors
\[
 \ell_i=(w_i\cdot w_i)^{-1}w_i^{\mathsf T},
 \qquad v_i=w_i,
\]
where the first expression is regarded as a \(1\times5\) row and
the second as a \(5\times1\) column. Then
\[
 \ell_i v_i=1,\qquad \ell_i v_j=0\quad(i\sim_Bj).
\]
Thus $h_3(B)\le5$.

\subsubsection{A rank-eleven, block-two representation of H}

The following eleven subsets are independent in \(H\), and every vertex
belongs to exactly two of them.
\begin{center}
\begin{tabular}{r@{\quad}l}
\toprule
\(a\)&Independent set \(I_a\)\\
\midrule
1&\{2,\,3,\,13,\,18,\,25\}\\
2&\{3,\,6,\,7,\,16,\,21\}\\
3&\{8,\,11,\,18,\,27,\,28\}\\
4&\{8,\,13,\,22,\,23,\,26\}\\
5&\{0,\,9,\,14,\,19,\,24\}\\
6&\{4,\,9,\,14,\,21,\,24\}\\
7&\{2,\,11,\,16,\,27,\,28\}\\
8&\{0,\,5,\,10,\,15,\,20\}\\
9&\{5,\,10,\,15,\,20,\,25\}\\
10&\{6,\,7,\,12,\,22,\,23\}\\
11&\{4,\,12,\,19,\,26\}\\
\bottomrule
\end{tabular}
\end{center}
Both claims can be checked directly from the adjacency table.
Let $e_1,\ldots,e_{11}$ be the standard coordinate vectors of $\F_3^{11}$.
For \(h\in V(H)\), let \(a(h)<b(h)\) be its two set indices and put
\[
 E_h=[\,e_{a(h)}\ \ e_{b(h)}\,]\in\F_3^{11\times2}.
\]
Then \(E_h^{\mathsf T}E_h=I_2\), and adjacent vertices have disjoint
index pairs, so \(E_h^{\mathsf T}E_{h'}=0\) on every \(H\)-edge.
Thus $h_3(H)\le11/2$; the same representation works over every field.

\subsection{Exact real Gram certificates}\label{subsix:sec:gram}

The following Gram-matrix criterion gives the theta bounds we need.
\begin{lemma}\label{subsix:lem:gram}
Suppose a real positive semidefinite matrix \(M\), indexed by \(V(X)\),
has diagonal entries \(t-1\) and entries \(-1\) on every edge of \(X\).
Then there are unit vectors \(u_v\) and a fixed unit vector \(c\),
called the handle, satisfying
\[
 u_v\perp u_w\quad(v\sim_Xw),\qquad
 \langle c,u_v\rangle=t^{-1/2}.
\]
Consequently $\overline\vartheta(X)\le t$ and $\C(X)\le t$.
\end{lemma}
\begin{proof}
Write \(M_{vw}=\langle z_v,z_w\rangle\), using a real Gram factorization.
Take \(u_v=(1,z_v)/\sqrt t\) and \(c=(1,0,\ldots,0)\).
The orthonormal-representation bound recalled in Supplementary
App.~\ref{sec:finiteproof} gives $\overline\vartheta(X)\le t$.
\end{proof}

Subsec.~\ref{subsix:app:matrices} supplies the following certificates.
\begin{center}
\begin{tabular}{c r r r c}
\toprule
Graph&Order&Scale $s_X$&Diagonal of $M_X$&Parameter $t_X$\\
\midrule
$B$&27&$24$&$104$&$16/3$\\
$D$&27&$100000000$&$419004112$&$b$\\
$H$&27&$100000000$&$419656120$&$q$\\
\bottomrule
\end{tabular}
\end{center}
Each $M_X$ is positive definite, has entry $-s_X$ on graph edges,
and has normalized diagonal $t_X-1$. The symmetry reductions below
verify positivity on blocks of size at most $6$ for $B,D$, and of sizes
$15$ and $12$ for $H$. Applying Lem.~\ref{subsix:lem:gram} to the
normalized matrices $M_X/s_X$ gives
\[
 \overline\vartheta(B)\le16/3,\qquad
 \overline\vartheta(D)\le b,\qquad
 \overline\vartheta(H)\le q.
\]
The maximum rule~\eqref{eq:max} therefore proves
\begin{equation}\label{subsix:eq:factoruppers}
 \C(G)\le\overline\vartheta(G)
 =\max\{\overline\vartheta(B),\overline\vartheta(D)\}\le16/3<6,
 \qquad \C(H)\le q<6.
\end{equation}

\subsection{The product bound}\label{subsix:sec:weighted}

The graph $G\times H$ is $(B\times H)\sqcup(D\times H)$.
Projecting to the appropriate factor in each component gives
\[
\begin{array}{c|cc}
 &h_3&\overline\vartheta\\ \hline
 B\times H&\le5&\le q\\
 D\times H&\le11/2&\le b .
\end{array}
\]
Apply~\eqref{eq:weighted-main} with these two spectral points. For every
$0\le\lambda\le1$,
\begin{equation}\label{subsix:eq:productupper}
 \C(G\times H)\le T_\lambda
 :=\max\{5^\lambda q^{1-\lambda},(11/2)^\lambda b^{1-\lambda}\}.
\end{equation}
This already controls arbitrary codes of every length, including words
whose coordinates lie in different components.

\subsection{Exact verification of the strict separation}

Choose the rational weight \(\lambda=13/1000\). The required strict
comparisons reduce to the rational inequalities
\begin{align}
 5^{13}\left(\frac{12991403}{2500000}\right)^{987}
 &<
 \left(\frac{5193957}{1000000}\right)^{1000},\label{subsix:eq:power1}\\
 \left(\frac{11}{2}\right)^{13}
 \left(\frac{32437757}{6250000}\right)^{987}
 &<
 \left(\frac{5193957}{1000000}\right)^{1000}.\label{subsix:eq:power2}
\end{align}
The inequalities~\eqref{subsix:eq:power1}--\eqref{subsix:eq:power2} are
checked by integer powers and cross multiplication; the exact
Python statements are given below. Taking positive \(1000\)-th roots
proves \(T_\lambda<L\). Also
\[
 5193957^2=26977189317849<27000000000000=27\cdot10^{12},
\]
so \(L<\sqrt{27}\). Combining \eqref{subsix:eq:Glower}, \eqref{subsix:eq:Hlower},
\eqref{subsix:eq:factoruppers}, and \eqref{subsix:eq:productupper} proves
Prop.~\ref{subsix:thm:main}.

To establish the sub-six spectral claim in \cref{sec:spectral-consequences}, choose the
single interpolating point $\psi\in\spec$ supplied by
\eqref{eq:interpolate} for $h_3,\overline\vartheta$ with weights
$13/1000,987/1000$. The component bounds above and the strict
comparisons~\eqref{subsix:eq:power1}--\eqref{subsix:eq:power2} give
$\psi(G\times H)\le T_{13/1000}<L$.
Every spectral point bounds capacity from above, so
$\psi(G),\psi(H)\ge\sqrt{27}$.
To bound their values from above, the identity fitting matrix gives
$h_3(D)\le27$, and the maximum rule gives $h_3(G)\le27$.
The existing certificates also give
$\overline\vartheta(G)\le16/3$, $h_3(H)\le11/2$ and
$\overline\vartheta(H)\le q<11/2$.
Thus the same interpolant satisfies
\[
 \psi(G)\le27^{13/1000}(16/3)^{987/1000}
 \le\frac{13}{1000}\,27+\frac{987}{1000}\,\frac{16}{3}
 =\frac{1123}{200}<6,
\]
by the weighted arithmetic--geometric mean inequality, while
\[
 \psi(H)\le(11/2)^{13/1000}q^{987/1000}<11/2<6.
\]
Thus one spectral point has product value below $\sqrt{27}$ and
both factor values in $[\sqrt{27},6)$, as asserted in
\cref{sec:spectral-consequences}.

\subsection{Compact integer matrix certificates}\label{subsix:app:matrices}

The scales and diagonals are those in Subsec.~\ref{subsix:sec:gram}.
Every edge entry is minus the scale. Only nonedge entries remain to be
specified. In each table, \texttt{i: j=value} gives an entry at an unordered
pair $\{i,j\}$; extend it to the whole orbit under the stated permutations,
and set $M_{ji}=M_{ij}$. For the group $\Gamma$ generated by those
permutations, this pair orbit is
$\{\{g(i),g(j)\}:g\in\Gamma\}$.
Each nonedge orbit occurs exactly once.

\subsubsection{The two source matrices}

Both $B$ and $D$ are preserved by the permutations
\begin{align*}
 a&=(4\,5)(6\,9)(7\,11)(8\,10)(17\,22)(18\,23)(19\,24)(20\,25)(21\,26)\\
 b_0&=(0\,17)(1\,19)(4\,23)(6\,25)(7\,12)(8\,26)(11\,16)(13\,24)(14\,22)\\
 c&=(0\,16)(2\,15)(4\,25)(5\,20)(6\,23)(7\,22)(9\,18)(11\,17)(12\,14).
\end{align*}
They generate a group $\Gamma$ of order $16$. The following $23$ and $19$
entries determine $M_B$ and $M_D$, respectively.

\noindent\textbf{Source certificate $M_B$:}
% BEGIN COMPACT B
\begin{Verbatim}[fontsize=\scriptsize,baselinestretch=1.03]
 0: 1=24 2=6 7=15 8=22 14=-9 16=12 18=18 19=27 20=-2
 1: 3=-31 4=4 19=11 21=5
 2: 3=-4 4=20 8=27
 3: 4=35
 4: 6=2 8=-20 9=19 23=23 26=34
 8: 10=6
\end{Verbatim}
% END COMPACT B
\noindent\textbf{Complement certificate $M_D$:}
% BEGIN COMPACT D
\begin{Verbatim}[fontsize=\scriptsize,baselinestretch=1.03]
 0: 3=92395858 4=145133660 6=-30688165 12=49988928 13=182327057 15=132541340
    17=84507591 21=76315484
 1: 2=3087050 6=138042021 8=109278319 13=-87574439
 2: 6=127958661 15=-3681718
 3: 8=40039826
 4: 5=47798171 10=83490663 25=123407703
 8: 26=-16181236
\end{Verbatim}
% END COMPACT D

Here is an elementary reduction of their positivity checks. Write
$A,B_0,C$ for the permutation matrices of $a,b_0,c$, and put $Z=(AB_0)^2$.
These are symmetric involutions; $C,Z$ commute with $A,B_0$, and
$AB_0=ZB_0A$. Thus $A$ and $B_0$ commute on the $Z=+I$ eigenspace
and anticommute on the $Z=-I$ eigenspace.

Let $\varepsilon,\delta,\alpha,\beta\in\{+1,-1\}$ label the eigenvalues
of $Z,C,A,B_0$, respectively. When $\varepsilon=+1$, the simultaneous
eigenspaces of these four commuting symmetric matrices form an orthogonal
decomposition. When $\varepsilon=-1$, decompose instead by the
eigenvalues $\delta,\alpha$ of $C,A$. The orthogonal map $B_0$ carries
each such space onto the space with signs $\delta,-\alpha$.
A $\Gamma$-invariant symmetric matrix commutes with every generator,
preserves these spaces, and has orthogonally equivalent restrictions
on each paired set of spaces. It therefore suffices to check one member
of each pair, explaining the multiplicity $2$ below.

To obtain explicit integer bases, apply
\[
 (I+\beta B_0)(I+\alpha A)(I+\delta C)(I+\varepsilon Z)
\]
to $e_0,e_1,\ldots,e_{26}$ when $\varepsilon=+1$, and omit the first
factor when $\varepsilon=-1$. Divide each nonzero column by the gcd of its
entries, then retain the first linearly independent columns. Denote the
resulting basis matrix by $P$. The nonzero spaces are listed below;
when $\varepsilon=-1$, it suffices to take $\alpha=+1$, with multiplicity
$2$. The dimensions, counted with multiplicity, sum to $27$.

For each certificate, $P^{\mathsf T}M_XP$ represents the restricted
quadratic form in the chosen basis, which need not be orthonormal.
Its positive definiteness is therefore equivalent to that of the
restriction. Divide this matrix by the positive gcd of its entries.
The last column records the leading principal minors, the determinants
of the upper-left square submatrices, for the resulting $M_B$ blocks;
all are positive. Applying the exact
elimination below to the $M_D$ blocks likewise gives only positive minors.
\begin{center}
\small
\begin{tabular}{ccccrrl}
\toprule
$\varepsilon$&$\delta$&$\alpha$&$\beta$&Order&Multiplicity&Minors for $M_B$\\
\midrule
$+$&$+$&$+$&$+$&6&1&$130,1596,5376,2192,96,112$\\
$+$&$+$&$+$&$-$&2&1&$101,13$\\
$+$&$+$&$-$&$+$&2&1&$115,66$\\
$+$&$-$&$+$&$+$&3&1&$29,28,20$\\
$+$&$-$&$+$&$-$&1&1&$1$\\
$+$&$-$&$-$&$+$&1&1&$1$\\
$-$&$+$&$+$&--&4&2&$149,7232,10816,3786$\\
$-$&$-$&$+$&--&2&2&$77,94$\\
\bottomrule
\end{tabular}
\end{center}
Thus both source certificates are positive definite, using blocks of
order at most six.

\Needspace{8\baselineskip}
\subsubsection{The partner matrix}

On $H$, let $\sigma$ fix $0,12,28$ and send every other vertex
$v$ to $29-v$. It preserves adjacency. The next $91$ entries, extended
under $\sigma$, specify $M_H$.
% BEGIN COMPACT H
\begin{Verbatim}[fontsize=\scriptsize,baselinestretch=1.03]
 0: 3=-92182906 5=112385240 9=105366413 10=260048876 11=-64530825
    14=103723987 28=24358075
 2: 3=66703246 11=207795723 13=104538630 14=-136454182 16=177512886
    18=115342751 19=-34383691 20=-100268450 21=76229467 24=-95624150
    25=120690256 26=152316144 27=18986268 28=103951980
 3: 4=-35794660 5=-98567795 6=149505933 7=190124261 10=-22994873 12=66710199
    13=124899718 16=61341650 18=120482232 19=-19609272 21=-23529434
    24=-98379162 25=42147361
 4: 9=272941041 11=23305245 12=-70035550 14=241110784 16=61213470
    19=135393685 21=30684152 24=262861878
 5: 8=35191781 10=321011184 14=-104169070 15=375111138 20=418775902
    23=-98757666 28=-47386121
 6: 7=381208558 11=-42891154 12=337627573 13=-146194020 14=-78782821
    16=19526989 19=-86283590 21=90671780 22=187815754 23=202130330
    28=-36203985
 7: 8=-59663804 9=-100785398 12=291996527 15=-44507231 16=55398098
    21=117706337 22=165147825
 8: 11=-56920029 13=259229426 14=-47902526 15=57463096 18=183702360
    19=-222235396 20=41480973 21=10556935 28=92833206
 9: 14=379995036 18=-98046773 19=316558361 28=-53434665
10: 12=-37639343 14=-98004443 15=294664038 16=-99343891 18=-90692258
    19=60606640
11: 16=300836730 18=35163766 28=237971338
13: 14=-57898358 28=23099334
\end{Verbatim}
% END COMPACT H

For the positivity check, order the vertices as $F,P,\sigma(P)$, where
\[
 F=(0,12,28),\qquad P=(2,3,4,5,6,7,8,9,10,11,13,14).
\]
The invariant matrix has the form
\[
 M_H=\begin{pmatrix}E&F_0&F_0\\
 F_0^{\mathsf T}&J&K\\F_0^{\mathsf T}&K&J\end{pmatrix}.
\]
Its quadratic form splits into symmetric and antisymmetric coordinates.
Consequently it is positive definite if and only if
\[
 S=\begin{pmatrix}2E&2F_0\\2F_0^{\mathsf T}&J+K\end{pmatrix}
 \quad\text{and}\quad T=J-K
\]
are positive definite. Indeed, for coordinates $(x,y,z)$ of dimensions
$3,12,12$, respectively,
\[
 \begin{pmatrix}x\\y\\z\end{pmatrix}^{\mathsf T}
 M_H\begin{pmatrix}x\\y\\z\end{pmatrix}
 =\frac12\begin{pmatrix}x\\y+z\end{pmatrix}^{\mathsf T}
 S\begin{pmatrix}x\\y+z\end{pmatrix}
 +\frac12(y-z)^{\mathsf T}T(y-z).
\]
The change of variables $(x,y,z)\mapsto(x,y+z,y-z)$ is invertible,
which proves the equivalence. Exact elimination gives $15$ and $12$ positive
leading principal minors, respectively, proving that $M_H$ is positive
definite.

\Needspace{10\baselineskip}
\subsection{Exact verification and reproducibility}\label{subsix:app:verification}

All entries are specified above. The permutations, orbit coverage,
graph edges and block dimensions are finite checks. For positive
definiteness, we use Bareiss's fraction-free elimination~\cite{Bareiss}.
Sylvester's determinant identity shows that, at stage $k$,
\texttt{pivot} is the leading principal minor of order $k+1$ and all
divisions are exact. Positivity of these minors proves positive
definiteness by Sylvester's criterion.
\begin{Verbatim}[fontsize=\footnotesize,baselinestretch=.96]
def positive_definite(M):
    n = len(M)
    assert all(len(row) == n for row in M)
    assert all(M[i][j] == M[j][i]
               for i in range(n) for j in range(n))
    T = [row[:] for row in M]
    previous = 1
    for k in range(n):
        pivot = T[k][k]
        assert pivot > 0
        for i in range(k + 1, n):
            for j in range(i, n):
                num = pivot*T[i][j] - T[i][k]*T[k][j]
                assert num % previous == 0
                T[i][j] = T[j][i] = num // previous
        previous = pivot
    return n
\end{Verbatim}
The following code checks the strict scalar comparisons.
\begin{Verbatim}[fontsize=\footnotesize,baselinestretch=.96]
from fractions import Fraction as F
b = F(32437757, 6250000)
q = F(12991403, 2500000)
L = F(5193957, 1000000)
assert F(5)**13 * q**987 < L**1000
assert F(11, 2)**13 * b**987 < L**1000
assert L*L < 27 and max(F(16, 3), b, q) < 6
\end{Verbatim}
The verification programs and machine-readable data are available in
the accompanying repository~\cite{ZeissData}.
The program \nolinkurl{verification/check_subsix.py} verifies all finite
inputs and the symmetry reductions. The independent
\nolinkurl{verification/audit_subsix_from_tex.py} reconstructs the printed
data and checks the full matrices by rational $LDL^{\mathsf T}$
factorization. Run either with Python 3.11 or later and assertions enabled:
do not pass \texttt{-O} or set \texttt{PYTHONOPTIMIZE}. Only the standard
library is required. Both checkers use exact arithmetic, without floating-point
operations or a numerical solver. They verify the matrix certificates
used by the analytical interpolation argument.


\paragraph{AI assistance and verification workflow.}
OpenAI Codex assisted with the development and review of the verification
programs and with exploration of candidate constructions and certificates.
The archived inputs are the explicit graph definitions, finite codes and
integer or rational certificate entries printed in these appendices.
The programs check these fixed inputs; they do not rely on an AI service,
a numerical solver or an AI-generated assessment of correctness at run time.
In particular, the source-derived checker reconstructs the printed matrices
independently of the modular checker. Integer and rational computations
establish the finite certificate inequalities. Decimal evaluation in the
uniform-threshold audit is diagnostic only. The spectral interpolation
and asymptotic arguments are proved analytically in the appendices and
are not certified by these finite programs.
